%%
%% This is file `sample-acmtog.tex',
%% generated with the docstrip utility.
%%
%% The original source files were:
%%
%% samples.dtx  (with options: `all,journal,bibtex,acmtog')
%% 
%% IMPORTANT NOTICE:
%% 
%% For the copyright see the source file.
%% 
%% Any modified versions of this file must be renamed
%% with new filenames distinct from sample-acmtog.tex.
%% 
%% For distribution of the original source see the terms
%% for copying and modification in the file samples.dtx.
%% 
%% This generated file may be distributed as long as the
%% original source files, as listed above, are part of the
%% same distribution. (The sources need not necessarily be
%% in the same archive or directory.)
%%
%%
%% Commands for TeXCount
%TC:macro \cite [option:text,text]
%TC:macro \citep [option:text,text]
%TC:macro \citet [option:text,text]
%TC:envir table 0 1
%TC:envir table* 0 1
%TC:envir tabular [ignore] word
%TC:envir displaymath 0 word
%TC:envir math 0 word
%TC:envir comment 0 0
%%
%% The first command in your LaTeX source must be the \documentclass
%% command.
%%
%% For submission and review of your manuscript please change the
%% command to \documentclass[manuscript, screen, review]{acmart}.
%%
%% When submitting camera ready or to TAPS, please change the command
%% to \documentclass[sigconf]{acmart} or whichever template is required
%% for your publication.
%%
%%
\documentclass[acmtog,nonacm]{acmart}
\usepackage{colortbl}
\usepackage{tikz}
\usepackage{multirow}
\usepackage{varwidth} %for the varwidth minipage environment
\usepackage{pifont} %for the varwidth minipage environment
\usepackage{boldline}
\usepackage{amsmath}
\usepackage{xcolor}
\usepackage[ruled]{algorithm2e} % For algorithms
\usepackage{wrapfig}
\usepackage{bm}
\usepackage{makecell}
\usepackage{arydshln}
\usepackage[utf8x]{inputenc}
\usepackage{siunitx}
\usepackage{afterpage}
\usepackage{graphicx}  % Add this to the preamble
\newcolumntype{P}[1]{>{\centering\arraybackslash}p{#1}}
%%
%% \BibTeX command to typeset BibTeX logo in the docs
\AtBeginDocument{%
  \providecommand\BibTeX{{%
    Bib\TeX}}}

\setlength{\textfloatsep}{4pt}
% \setlength{\abovedisplayskip}{0pt} % Space above equations
% \setlength{\belowdisplayskip}{0pt} % Space below equations

\graphicspath{ {../} }
\captionsetup[figure]{aboveskip=4pt, belowskip=2pt}

%% Rights management information.  This information is sent to you
%% when you complete the rights form.  These commands have SAMPLE
%% values in them; it is your responsibility as an author to replace
%% the commands and values with those provided to you when you
%% complete the rights form.
\setcopyright{acmlicensed}
\acmJournal{TOG}
\acmYear{2025} \acmVolume{44} \acmNumber{4} \acmArticle{} \acmMonth{8}\acmDOI{10.1145/3731163}

%%
%% These commands are for a JOURNAL article.
\acmJournal{TOG}
\acmVolume{44}
\acmNumber{4}
\acmArticle{111}
\acmMonth{8}

%%
%% Submission ID.
%% Use this when submitting an article to a sponsored event. You'll
%% receive a unique submission ID from the organizers
%% of the event, and this ID should be used as the parameter to this command.

\acmSubmissionID{525}
%%\acmSubmissionID{123-A56-BU3}

%%
%% For managing citations, it is recommended to use bibliography
%% files in BibTeX format.
%%
%% You can then either use BibTeX with the ACM-Reference-Format style,
%% or BibLaTeX with the acmnumeric or acmauthoryear sytles, that include
%% support for advanced citation of software artefact from the
%% biblatex-software package, also separately available on CTAN.
%%
%% Look at the sample-*-biblatex.tex files for templates showcasing
%% the biblatex styles.
%%

%%
%% The majority of ACM publications use numbered citations and
%% references.  The command \citestyle{authoryear} switches to the
%% "author year" style.
%%
%% If you are preparing content for an event
%% sponsored by ACM SIGGRAPH, you must use the "author year" style of
%% citations and references.
\citestyle{acmauthoryear}


%%
%% end of the preamble, start of the body of the document source.
\begin{document}

%%
%% The "title" command has an optional parameter,
%% allowing the author to define a "short title" to be used in page headers.
%\title{Gaussian wave splatting: \\Computer-generated holography using Gaussian splats}
%\title{HoloSplats: Gaussian Wave Splatting for Computer-generated Holography}
\title{Gaussian Wave Splatting for Computer-Generated Holography}

%%
%% The "author" command and its associated commands are used to define
%% the authors and their affiliations.
%% Of note is the shared affiliation of the first two authors, and the
%% "authornote" and "authornotemark" commands
%% used to denote shared contribution to the research.



%%% Standard 
\author{Suyeon Choi}
\email{suyeon@stanford.edu}
\authornote{denotes equal contribution.}
\orcid{0000-0001-9030-0960}
\affiliation{
  \institution{Stanford University}
  \country{USA}
}
\author{Brian Chao}
\email{brianchc@stanford.edu}
\orcid{0000-0002-4581-6850}
\authornotemark[1]
\affiliation{
  \institution{Stanford University}
  \country{USA}
}
\author{Jacqueline Yang}
\email{jyang01@stanford.edu}
\orcid{0009-0002-3101-3026}
\affiliation{
  \institution{Stanford University}
  \country{USA}
}
\author{Manu Gopakumar}
\email{manugopa@stanford.edu}
\orcid{0000-0001-9017-4968}
\affiliation{
  \institution{Stanford University}
  \country{USA}
}
\author{Gordon Wetzstein}
\email{gordon.wetzstein@stanford.edu}
\orcid{0000-0002-9243-6885}
\affiliation{
  \institution{Stanford University}
  \country{USA}
}



% \author{Suyeon Choi}
% \email{suyeon@stanford.edu}
% \authornote{denotes equal contribution.}
% \orcid{0000-0001-9030-0960}
% \author{Brian Chao}
% \email{brianchc@stanford.edu}
% \orcid{0000-0002-4581-6850}
% \authornotemark[1]
% \author{Jacqueline Yang}
% \email{jyang01@stanford.edu}
% \orcid{0009-0002-3101-3026}
% \author{Manu Gopakumar}
% \email{manugopa@stanford.edu}
% \orcid{0000-0001-9017-4968}
% \author{Gordon Wetzstein}
% \email{gordon.wetzstein@stanford.edu}
% \orcid{0000-0002-9243-6885}
% \affiliation{
%   \institution{Stanford University}
%   \country{USA}
% }





%%
%% By default, the full list of authors will be used in the page
%% headers. Often, this list is too long, and will overlap
%% other information printed in the page headers. This command allows
%% the author to define a more concise list
%% of authors' names for this purpose.
\renewcommand{\shortauthors}{Choi, Chao et al.}

%%
%% The abstract is a short summary of the work to be presented in the
%% article.
\begin{abstract}
State-of-the-art neural rendering methods optimize Gaussian scene representations from a few photographs for novel-view synthesis. Building on these representations, we develop an efficient algorithm, dubbed Gaussian Wave Splatting, to turn these Gaussians into holograms. Unlike existing computer-generated holography (CGH) algorithms, Gaussian Wave Splatting supports accurate occlusions and view-dependent effects for photorealistic scenes by leveraging recent advances in neural rendering.  
%
Specifically, we derive a closed-form solution for a 2D Gaussian-to-hologram transform that supports occlusions and alpha blending. Inspired by classic computer graphics techniques, we also derive an efficient approximation of the aforementioned process in the Fourier domain that is easily parallelizable and implement it using custom CUDA kernels. 
%
By integrating emerging neural rendering pipelines with holographic display technology, our Gaussian-based CGH framework paves the way for next-generation holographic displays.
%
%We aim to develop efficient and high-quality rendering algorithms for holographic displays. Current computer-generated holography (CGH) pipelines capable of providing high-quality view-dependent effects are computationally slow and rely on limited data formats, such as dense light fields, which are not readily available. While recent advancements in deep-learning-based CGH methods have shown promise in terms of efficiency, these approaches rely on RGBD data, focal stack, or point clouds, offering only partial support for accurate occlusion and directional effects. Our approach is inspired by and leverages recent advances in neural rendering to address these limitations. Specifically, we derive an analytical solution for CGH using anisotropic Gaussian primitives as input, and implement it using a custom CUDA kernel to rasterize each pixel in the Fourier domain efficiently. Starting with captured photographs of real or rendered scenes, our method transforms them into a high-quality 3D Gaussian representation, and performs the developed CGH algorithms to convert this representation into a set of holograms. This novel pipeline bridges the gap between computational efficiency, data bandwidth, and visual fidelity, enabling photorealistic computer-generated holography for captured scenes. By integrating emerging neural rendering techniques with holographic display technology, our method not only overcomes the shortcomings of existing approaches but also paves the way for practical, high-quality holographic rendering.
\end{abstract}

%%
%% The code below is generated by the tool at http://dl.acm.org/ccs.cfm.
%% Please copy and paste the code instead of the example below.
%%
\begin{CCSXML}
<ccs2012>
   <concept>
       <concept_id>10010147.10010371</concept_id>
       <concept_desc>Computing methodologies~Computer graphics</concept_desc>
       <concept_significance>500</concept_significance>
       </concept>
   <concept>
       <concept_id>10010583.10010786</concept_id>
       <concept_desc>Hardware~Emerging technologies</concept_desc>
       <concept_significance>300</concept_significance>
       </concept>
 </ccs2012>
\end{CCSXML}



%%%%%%%% Original %%%%%%%%%
\ccsdesc[500]{Computing methodologies~Computer graphics}
\ccsdesc[300]{Hardware~Emerging technologies}

%%
%% Keywords. The author(s) should pick words that accurately describe
%% the work being presented. Separate the keywords with commas.
\keywords{computational displays, holography, virtual reality, augmented reality, Gaussian splatting, neural rendering}



% %%%%%%%%% To see if we can save space %%%%%%%%
% \ccsdesc[300]{Hardware~Emerging technologies}
% % \keywords{holography, neural rendering}





\received{23 January 2025}
\received[revised]{12 March 2025}
\received[accepted]{5 June 2025}

\begin{teaserfigure}
  \centering
	\includegraphics[width=\columnwidth]{main_paper/figures_camera_ready/teaser.pdf}
   \caption{
   Emerging Gaussian splatting techniques take a few casually captured photographs as input and output a Gaussian scene representation (left). Our Gaussian Wave Splatting approach turns these types of representations into holograms that can be directly displayed on emerging holographic displays (center). The color bar below the hologram denotes phase values of the phase-only hologram. Our Gaussian-based computer-generated holography algorithm reconstructs photorealistic 3D scenes with refocusing capabilities as seen in the experimentally captured results on the right. }%\textcolor{purple}{GW: change ``Captured 3D Focal Stack'' in lower right label to ``Captured 3D Hologram''.
   %
%   Gaussian Wave Splatting (GWS) takes an optimized 2DGS model as input and directly converts it to 3D holograms. GWS works with an explicit 2D Gaussians primitive-based representation unlike other neural networks-based scene representations, thus no intermediate renders such as RGBD images or light fields are required to generate 3D holograms. GWS-generated 3D holograms reconstructs sharp details in in-focus regions and synthesizes more accurate blur in defocus regions compared to prior primitives-based CGH methods. 
% }
  \label{fig:teaser}
\end{teaserfigure} 

%%
%% This command processes the author and affiliation and title
%% information and builds the first part of the formatted document.
\maketitle
\input{definitions}

\section{Introduction}
\input{main_paper/sections/1_intro}

\section{Related Work}
\input{main_paper/sections/2_previous_work}

\section{Background on Gaussian Splatting}
\input{main_paper/sections/3_gs}

\section{Wave Splatting}
\input{main_paper/sections/4_method_rewritten}

% \section{Wave Splatting}
% \input{main_paper/sections/4_method}

%\vspace{-3pt}
\section{Experiments}
\label{sec:results}
\input{main_paper/sections/5_results}

%\vspace{-3pt}
\section{Discussion}
\input{main_paper/sections/6_discussion}

% \begin{description}
% \item[\texttt{sidebar}:]  Place formatted text in the margin.
% \item[\texttt{marginfigure}:] Place a figure in the margin.
% \item[\texttt{margintable}:] Place a table in the margin.
% \end{description}
%%
%% The next two lines define the bibliography style to be used, and
%% the bibliography file.
\begin{acks}

% We thank Seung-Woo Nam and Gun-Yeal Lee for fruitful discussions. Brian Chao is supported by a Stanford Graduate Fellowship and an NSF GRFP. Jacqueline Yang is also supported by a Stanford Graduate Fellowship and the NSF GRFP.

We thank Seung-Woo Nam and Gun-Yeal Lee for fruitful discussions. Brian Chao and Jacqueline Yang are supported by Stanford Graduate Fellowships and the NSF Graduate Research Fellowship Program (GRFP). We also thank EssilorLuxottica and Meta for their support.



\end{acks}

\bibliographystyle{ACM-Reference-Format}
% \bibliography{sample-base}
\bibliography{bibs}

%%%%%%%%%%%%%%%%%%%%%%%%%%%%%%%% Expected to be in figure-only pages %%%%%%%%%%%%%%%%%%%%%%%%%%%%%%%%%%
%%
%% The acknowledgments section is defined using the "acks" environment
%% (and NOT an unnumbered section). This ensures the proper
%% identification of the section in the article metadata, and the
%% consistent spelling of the heading.






%%
%% If your work has an appendix, this is the place to put it.
% \appendix

% \section{Research Methods}
% \clearpage
% \afterpage{%
\begin{figure*}[ht!]
    \centering
    % \includegraphics[width=0.92\textwidth]{main_paper/figures/main_baseline_captured.pdf}
    \includegraphics[width=1.0\textwidth]{main_paper/figures_camera_ready/fig5_results.pdf}
    \caption{\textbf{Experimentally captured focal stacks of holograms generated using different primitive-based CGH algorithms. } For each CGH baseline, we show the target rendering of its input 3D representation using a pinhole camera model (\textit{top row}). Sparse point-cloud representations used in point-based methods~\cite{chen2009computer} inevitably lead to holes in the rendered target image. Polygon-based methods~\cite{matsushima2009extremely} produce overly smooth results unless an excessive number of polygons are used. Our Gaussian Wave Splatting methods achieve the best results with sharp in-focus regions and more accurate blur in defocus regions. Please refer to the supplementary materials for extensive baseline comparisons.}
    \label{fig:baseline_simulation}
\end{figure*}

\begin{figure*}[ht!]
    \centering
    \includegraphics[width=1.0\textwidth]{main_paper/figures_camera_ready/fig6_nvs.pdf}
    \caption{\textbf{Experimentally captured focal stacks of exact GWS holograms from different novel viewpoints.} We generate holograms from \revision{multiple} novel viewpoints \revision{using our exact GWS method applied to a single scene representation} and capture the corresponding 3D focal stacks at each viewpoint. \revision{The results show that }GWS achieves accurate 3D refocusing effects \revision{across} all viewpoints, demonstrating the robustness of our method. \revision{See the supplementary video for an additional visualization.}}%\sy{For revision, let's rotate the text by 90 degrees, and I would vote for unbold "captured" for camera-ready version.} }    
    \label{fig:nvs_fs_captured}
\end{figure*}

% \clearpage
% }
\end{document}
\endinput
%%
%% End of file `sample-acmtog.tex'.
